\section{Background}

\subsection{Cardinality estimation and native statistics}

The stock PostgreSQL planner combines relation-level information, per-column
statistics, and multivariate statistics to estimate predicate selectivities.
When predicates refer to correlated columns, multiplying independent
selectivities can be a poor approximation. PostgreSQL's extended statistics
objects provide a native representation for selected column groups; the
catalog definition and the data produced by \texttt{ANALYZE} are separate
pieces of state \cite{PostgreSQL16Docs,PostgreSQL16Catalog}. The system under
study uses the MCV and functional-dependency kinds supported by the patched
what-if path.

Creating a statistics object is not the same as knowing that it is useful.
The object has to cover workload predicates, be built at an adequate target,
and interact with the planner's selection and combination rules. Enumerating
all column groups is infeasible for realistic schemas, so the advisor freezes a
candidate universe and treats selection as a bounded design problem.

\subsection{What-if design and operational separation}

Hypothetical indexes have long illustrated the value of asking a planner what
it would do if a physical object existed, without paying the full build cost;
HypoPG is a PostgreSQL example \cite{HypoPG}. The present system applies a
related idea to native extended-statistics payloads. Its hypothetical objects
are backend-local and are consumed by the same planner selectivity paths,
while deployment remains ordinary PostgreSQL DDL.

The separation is important for operational reasons. The optimizer explores
on a controlled sample and workload snapshot; deployment operates on the
production population and can be reviewed by a DBA. The deployment contract
is add-only with respect to advisor ownership. In particular, an arbitrary
existing production statistics object does not make preflight fail, but a
deterministic-name collision with a current recommendation is fail-closed.

\subsection{Research questions}

We organize the study around five questions:

\begin{description}
  \item[RQ1 (effectiveness)] How much cardinality-estimation error can
  workload-optimized native extended statistics eliminate?
  \item[RQ2 (transfer)] Do designs optimized on a fixed sample retain their
  benefits when native payloads are independently recomputed on full data?
  \item[RQ3 (fidelity)] How faithfully does catalogless hypothetical evaluation
  reproduce physical PostgreSQL statistics behavior when data, payload
  semantics, ordinary statistics state, and design are controlled?
  \item[RQ4 (necessity)] What does planner-in-the-loop search contribute beyond
  simple statistics-selection heuristics?
  \item[RQ5 (practicality)] What operational cost is paid for the obtained
  improvement?
\end{description}
